\SourceLeaf{442}%
then we have
\[
\begin{aligned}
\chi_1&=\chi+\operatorname{Area}L'L'_1P'_1P'\\
\psi'-\varphi'+A&=LL'_1L'L+\chi_1-\chi;
\end{aligned}
\]
but
\[
\begin{aligned}
\varphi'&=\chi+B\\
\psi'&=\chi_1+B_1,
\end{aligned}
\]
hence
\[
B_1-B+A=LL'_1L'L.
\]
The angles of the triangle $LL'L'_1$ are evidently
\[
A,\qquad 180^{0}-B,\qquad B_1,
\]
so their sum is
\[
180^{0}+LL'_1L'L.
\]

The form of the proof will require some modification and explanation if the point $(3)$ lies inside the triangle. In general, however, we conclude:
\begin{quote}
»The sum of the three angles of a triangle formed by shortest lines on any curved surface is equal to the sum of $180^{0}$ and the area of the triangle on the auxiliary sphere. The boundary of the latter triangle is formed by the points $L$ corresponding to the points on the boundary of the first triangle. The area in question is to be regarded as positive or negative according as its boundary is traversed in the same sense as the figure or in the opposite sense«.
\end{quote}
\noindent\hspace*{1.2em}It also follows readily from this that the sum of all the angles of a polygon of $n$ sides that are shortest lines on the curved surface [equals] the sum of $(n-2)180^{0}+$ the area of the polygon on the sphere, etc.

\Needspace{5\baselineskip}
\begin{center}16.\end{center}
\noindent\hspace*{1.2em}If one curved surface can be developed completely onto another, then clearly all lines on the first surface preserve their lengths when transferred to the other, as do the angles formed by the
